\begin{figure}[h]
\centering
\includegraphics[width=0.7\columnwidth]{figures/fig_001.png}
\caption{fig\_001: Grouped bar chart comparing LLM performance on MATH-500 and AIME24, used as the source figure for the example capability questions in Table~\ref{tab:capability-categories}.}
\label{fig:capability-example}
\end{figure}

\begin{table*}[h]
\centering
\small
\begin{tabular}{@{}p{2.2cm}p{4.8cm}p{7cm}@{}}
\toprule
\textbf{Category} & \textbf{Description} & \textbf{Example (Figure~\ref{fig:capability-example})} \\
\midrule
Counting & Enumerate discrete visual elements such as bars, line series, segments, or legend entries. & \emph{How many models in the figure are within 1\% of another model in the same subplot?} \\
\addlinespace
Computation & Perform arithmetic on values read from the chart (differences, ratios, sums). & \emph{What is the percentage difference between the highest-performing and lowest-performing model in the MATH-500 dataset?} \\
\addlinespace
Comparison & Make relative judgements between chart elements (which is larger, rank ordering). & \emph{Which pair of models has the smallest gap in performance in chart (e), and what is the difference?} \\
\addlinespace
Pattern analysis & Identify trends, inflection points, or anomalies across the chart. & \emph{In the AIME24 subplots, do models that have similar scores in one subplot also remain close in the other?} \\
\bottomrule
\end{tabular}
\caption{Capability question categories with example questions drawn from fig\_001, a grouped bar chart comparing LLM performance on MATH-500 and AIME24.}
\label{tab:capability-categories}
\end{table*}
